\newpage
\section{Scriptural Date and Location}
\zlabel{triptych:concise:chronology:start}
% Generated. Do not edit.
% Deterministic projection of research/chronology.toml.
% schema: 5
% document: liturgy/roman-rite/postconciliar/roman-missal-third-edition-en-us-2011/propers/general-calendar/pc-r07-all-saints-abc
% generated-by: tools/tpt proper-chronology annotations
\newcommand{\chronologyannotationclaim}[6]{#6}
\newcommand{\chronologyannotationreach}[3]{}
\newcommand{\chronologyannotationgroup}[3]{#3}
\newcommand{\chronologyannotation}[1]{%
  \ifcsname triptychchronologyannotation@#1\endcsname
    \csname triptychchronologyannotation@#1\endcsname
  \else
    \PackageError{triptych}{No chronology annotation for #1}{}%
  \fi}
% appointment-dependency: guidance/liturgy/postconciliar-propers-registry.md sha256=ffe8a6d8b7d79c652bdc49b3d5f6d4dbb32e3506abc7ab0310e745c2b7306ad4
% appointment-dependency: src/gpt/liturgy/roman-rite/postconciliar/roman-missal-third-edition-en-us-2011/propers/general-calendar/pc-r07-all-saints-abc/instance/manifest.md sha256=9ae6055704d1fe2aa3949d5e91e816673c113bc775b3a39b1bc145c4290b6630
% appointment-dependency: src/gpt/liturgy/roman-rite/postconciliar/roman-missal-third-edition-en-us-2011/propers/general-calendar/pc-r07-all-saints-abc/propers/verified.md sha256=574617c10630e751cb43a297cfec38b5c622631c3dd5ca4b8092d3f70c5441b5
% appointment-dependency: src/gpt/liturgy/roman-rite/postconciliar/roman-missal-third-edition-en-us-2011/propers/general-calendar/pc-r07-all-saints-abc/research/chronology-inputs.toml sha256=f30b32a467cb7d56012eddb3291e509def6fa3b786ee844bead09d3b8452e975
% appointment-dependency: src/gpt/liturgy/roman-rite/postconciliar/roman-missal-third-edition-en-us-2011/propers/general-calendar/shared/formularies/pc-r07-all-saints/propers/verified.md sha256=3025d42355a3162576f155c12ed1d0eadca32ec592546ef8ca6bf7c7066e5630
% appointment-dependency: src/gpt/liturgy/roman-rite/postconciliar/roman-missal-third-edition-en-us-2011/propers/registry/README.md sha256=4e2d92326dece103e1475e70ffb5e447815f9b12d697926e48a54850a1641208
% appointment-dependency: src/gpt/liturgy/roman-rite/postconciliar/roman-missal-third-edition-en-us-2011/propers/registry/formula-dispositions.md sha256=de58754f5a8ab61dffa3f4a6a479c843fe3763b6a5b37099282960de71d93989
% appointment-note: entrance-antiphon: Composed Missal text without an appointed biblical citation; theological echoes and compiler chant directions do not create appointments.
% appointment-note: collect: Composed Missal text without an appointed biblical citation; theological echoes and compiler chant directions do not create appointments.
% appointment-note: first-reading: Revelation 7:2-4,9-14; appointed; vulgate numbering; canonical verse queries, not appointed wording.
% appointment-note: responsorial-psalm: Psalm 24:1bc-2,3-4ab,5-6; appointed; hebrew numbering; whole-verse query envelope, not subverse chronology or appointed wording.
% appointment-note: responsorial-psalm: Cf. Psalm 24:6; adaptation; hebrew numbering; canonical verse queries, not appointed wording.
% appointment-note: second-reading: 1 John 3:1-3; appointed; vulgate numbering; canonical verse queries, not appointed wording.
% appointment-note: gospel-acclamation: Matthew 11:28; appointed; vulgate numbering; canonical verse queries, not appointed wording.
% appointment-note: gospel: Matthew 5:1-12a; appointed; vulgate numbering; whole-verse query envelope, not subverse chronology or appointed wording.
% appointment-note: prayer-over-offerings: Composed Missal text without an appointed biblical citation; theological echoes and compiler chant directions do not create appointments.
% appointment-note: proper-preface: Composed Missal text without an appointed biblical citation; theological echoes and compiler chant directions do not create appointments.
% appointment-note: communion-antiphon: Matthew 5:8-10; appointed; vulgate numbering; canonical verse queries, not appointed wording.
% appointment-note: prayer-after-communion: Composed Missal text without an appointed biblical citation; theological echoes and compiler chant directions do not create appointments.
% appointment-note: solemn-blessing: Composed Missal text without an appointed biblical citation; theological echoes and compiler chant directions do not create appointments.
% first-reading: First Reading
\expandafter\def\csname triptychchronologyannotation@first-reading\endcsname{%
  \chronologyannotationgroup{composition}{disputed}{\textbf{Composition} -- disputed: \chronologyannotationreach{Apoc.7.10}{Apoc}{true}\chronologyannotationreach{Apoc.7.11}{Apoc}{true}\chronologyannotationreach{Apoc.7.12}{Apoc}{true}\chronologyannotationreach{Apoc.7.13}{Apoc}{true}\chronologyannotationreach{Apoc.7.14}{Apoc}{true}\chronologyannotationreach{Apoc.7.2}{Apoc}{true}\chronologyannotationreach{Apoc.7.3}{Apoc}{true}\chronologyannotationreach{Apoc.7.4}{Apoc}{true}\chronologyannotationreach{Apoc.7.9}{Apoc}{true}\chronologyannotationclaim{composition.apocalypse}{composition}{catholic-traditional-v1}{disputed}{the reign of Domitian ... A.D. 81-96}{A.D. 81--96}; \chronologyannotationreach{Apoc.7.10}{Apoc}{true}\chronologyannotationreach{Apoc.7.11}{Apoc}{true}\chronologyannotationreach{Apoc.7.12}{Apoc}{true}\chronologyannotationreach{Apoc.7.13}{Apoc}{true}\chronologyannotationreach{Apoc.7.14}{Apoc}{true}\chronologyannotationreach{Apoc.7.2}{Apoc}{true}\chronologyannotationreach{Apoc.7.3}{Apoc}{true}\chronologyannotationreach{Apoc.7.4}{Apoc}{true}\chronologyannotationreach{Apoc.7.9}{Apoc}{true}\chronologyannotationclaim{composition.apocalypse}{composition}{catholic-traditional-v1}{disputed}{from the year 90 to 100 (approximately)}{c. A.D. 90--100}.}%
}
% responsorial-psalm: Responsorial Psalm
\expandafter\def\csname triptychchronologyannotation@responsorial-psalm\endcsname{%
  \chronologyannotationgroup{historical-setting}{research-pending}{\textbf{Historical setting} -- Occasion date unresolved.}%
  \space \chronologyannotationgroup{composition}{preferred}{\textbf{Composition}: \chronologyannotationreach{Ps.23.1}{Ps}{true}\chronologyannotationreach{Ps.23.2}{Ps}{true}\chronologyannotationreach{Ps.23.3}{Ps}{true}\chronologyannotationreach{Ps.23.4}{Ps}{true}\chronologyannotationreach{Ps.23.5}{Ps}{true}\chronologyannotationreach{Ps.23.6}{Ps}{true}\chronologyannotationclaim{critical.psalms.latest-composition-boundary}{composition}{catholic-critical-v1}{preferred}{before the Maccabean period, around 165 B.C.; no individual psalm can be dated securely}{Before c. 165 B.C.}}%
}
% second-reading: Second Reading
\expandafter\def\csname triptychchronologyannotation@second-reading\endcsname{%
  \chronologyannotationgroup{composition}{preferred}{\textbf{Composition}: \chronologyannotationreach{1John.3.1}{1John}{true}\chronologyannotationreach{1John.3.2}{1John}{true}\chronologyannotationreach{1John.3.3}{1John}{true}\chronologyannotationclaim{composition.first-epistle-of-st-john}{composition}{catholic-traditional-v1}{preferred}{from the year 90 to 100 (approximately)}{c. A.D. 90--100}.}%
}
% gospel-acclamation: Gospel acclamation
\expandafter\def\csname triptychchronologyannotation@gospel-acclamation\endcsname{%
  \chronologyannotationgroup{composition}{disputed}{\textbf{Composition} -- disputed: \chronologyannotationreach{Matt.11.28}{Matt}{true}\chronologyannotationclaim{composition.gospel-of-matthew}{composition}{catholic-traditional-v1}{disputed}{about A.D. 38-45}{c. A.D. 38--45}; \chronologyannotationreach{Matt.11.28}{Matt}{true}\chronologyannotationclaim{composition.gospel-of-matthew}{composition}{catholic-traditional-v1}{disputed}{about the year 40-42}{c. A.D. 40--42}; \chronologyannotationreach{Matt.11.28}{Matt}{true}\chronologyannotationclaim{composition.gospel-of-matthew}{composition}{catholic-traditional-v1}{disputed}{the years 40-45}{A.D. 40--45}; \chronologyannotationreach{Matt.11.28}{Matt}{true}\chronologyannotationclaim{composition.gospel-of-matthew}{composition}{catholic-traditional-v1}{disputed}{about the year 60-68}{c. A.D. 60--68}; \chronologyannotationreach{Matt.11.28}{Matt}{true}\chronologyannotationclaim{composition.gospel-of-matthew}{composition}{catholic-traditional-v1}{disputed}{about the years 64-67}{c. A.D. 64--67}; \chronologyannotationreach{Matt.11.28}{Matt}{true}\chronologyannotationclaim{composition.gospel-of-matthew}{composition}{catholic-traditional-v1}{disputed}{about the year 50}{c. A.D. 50}.}%
}
% gospel: Gospel
\expandafter\def\csname triptychchronologyannotation@gospel\endcsname{%
  \chronologyannotationgroup{narrated-event}{research-pending}{\textbf{Event} -- No narrated-event date in the chronology corpus.}%
  \space \chronologyannotationgroup{composition}{disputed}{\textbf{Composition} -- disputed: \chronologyannotationreach{Matt.5.1}{Matt}{true}\chronologyannotationreach{Matt.5.10}{Matt}{true}\chronologyannotationreach{Matt.5.11}{Matt}{true}\chronologyannotationreach{Matt.5.12}{Matt}{true}\chronologyannotationreach{Matt.5.2}{Matt}{true}\chronologyannotationreach{Matt.5.3}{Matt}{true}\chronologyannotationreach{Matt.5.4}{Matt}{true}\chronologyannotationreach{Matt.5.5}{Matt}{true}\chronologyannotationreach{Matt.5.6}{Matt}{true}\chronologyannotationreach{Matt.5.7}{Matt}{true}\chronologyannotationreach{Matt.5.8}{Matt}{true}\chronologyannotationreach{Matt.5.9}{Matt}{true}\chronologyannotationclaim{composition.gospel-of-matthew}{composition}{catholic-traditional-v1}{disputed}{about A.D. 38-45}{c. A.D. 38--45}; \chronologyannotationreach{Matt.5.1}{Matt}{true}\chronologyannotationreach{Matt.5.10}{Matt}{true}\chronologyannotationreach{Matt.5.11}{Matt}{true}\chronologyannotationreach{Matt.5.12}{Matt}{true}\chronologyannotationreach{Matt.5.2}{Matt}{true}\chronologyannotationreach{Matt.5.3}{Matt}{true}\chronologyannotationreach{Matt.5.4}{Matt}{true}\chronologyannotationreach{Matt.5.5}{Matt}{true}\chronologyannotationreach{Matt.5.6}{Matt}{true}\chronologyannotationreach{Matt.5.7}{Matt}{true}\chronologyannotationreach{Matt.5.8}{Matt}{true}\chronologyannotationreach{Matt.5.9}{Matt}{true}\chronologyannotationclaim{composition.gospel-of-matthew}{composition}{catholic-traditional-v1}{disputed}{about the year 40-42}{c. A.D. 40--42}; \chronologyannotationreach{Matt.5.1}{Matt}{true}\chronologyannotationreach{Matt.5.10}{Matt}{true}\chronologyannotationreach{Matt.5.11}{Matt}{true}\chronologyannotationreach{Matt.5.12}{Matt}{true}\chronologyannotationreach{Matt.5.2}{Matt}{true}\chronologyannotationreach{Matt.5.3}{Matt}{true}\chronologyannotationreach{Matt.5.4}{Matt}{true}\chronologyannotationreach{Matt.5.5}{Matt}{true}\chronologyannotationreach{Matt.5.6}{Matt}{true}\chronologyannotationreach{Matt.5.7}{Matt}{true}\chronologyannotationreach{Matt.5.8}{Matt}{true}\chronologyannotationreach{Matt.5.9}{Matt}{true}\chronologyannotationclaim{composition.gospel-of-matthew}{composition}{catholic-traditional-v1}{disputed}{the years 40-45}{A.D. 40--45}; \chronologyannotationreach{Matt.5.1}{Matt}{true}\chronologyannotationreach{Matt.5.10}{Matt}{true}\chronologyannotationreach{Matt.5.11}{Matt}{true}\chronologyannotationreach{Matt.5.12}{Matt}{true}\chronologyannotationreach{Matt.5.2}{Matt}{true}\chronologyannotationreach{Matt.5.3}{Matt}{true}\chronologyannotationreach{Matt.5.4}{Matt}{true}\chronologyannotationreach{Matt.5.5}{Matt}{true}\chronologyannotationreach{Matt.5.6}{Matt}{true}\chronologyannotationreach{Matt.5.7}{Matt}{true}\chronologyannotationreach{Matt.5.8}{Matt}{true}\chronologyannotationreach{Matt.5.9}{Matt}{true}\chronologyannotationclaim{composition.gospel-of-matthew}{composition}{catholic-traditional-v1}{disputed}{about the year 60-68}{c. A.D. 60--68}; \chronologyannotationreach{Matt.5.1}{Matt}{true}\chronologyannotationreach{Matt.5.10}{Matt}{true}\chronologyannotationreach{Matt.5.11}{Matt}{true}\chronologyannotationreach{Matt.5.12}{Matt}{true}\chronologyannotationreach{Matt.5.2}{Matt}{true}\chronologyannotationreach{Matt.5.3}{Matt}{true}\chronologyannotationreach{Matt.5.4}{Matt}{true}\chronologyannotationreach{Matt.5.5}{Matt}{true}\chronologyannotationreach{Matt.5.6}{Matt}{true}\chronologyannotationreach{Matt.5.7}{Matt}{true}\chronologyannotationreach{Matt.5.8}{Matt}{true}\chronologyannotationreach{Matt.5.9}{Matt}{true}\chronologyannotationclaim{composition.gospel-of-matthew}{composition}{catholic-traditional-v1}{disputed}{about the years 64-67}{c. A.D. 64--67}; \chronologyannotationreach{Matt.5.1}{Matt}{true}\chronologyannotationreach{Matt.5.10}{Matt}{true}\chronologyannotationreach{Matt.5.11}{Matt}{true}\chronologyannotationreach{Matt.5.12}{Matt}{true}\chronologyannotationreach{Matt.5.2}{Matt}{true}\chronologyannotationreach{Matt.5.3}{Matt}{true}\chronologyannotationreach{Matt.5.4}{Matt}{true}\chronologyannotationreach{Matt.5.5}{Matt}{true}\chronologyannotationreach{Matt.5.6}{Matt}{true}\chronologyannotationreach{Matt.5.7}{Matt}{true}\chronologyannotationreach{Matt.5.8}{Matt}{true}\chronologyannotationreach{Matt.5.9}{Matt}{true}\chronologyannotationclaim{composition.gospel-of-matthew}{composition}{catholic-traditional-v1}{disputed}{about the year 50}{c. A.D. 50}.}%
}
% communion-antiphon: Communion antiphon
\expandafter\def\csname triptychchronologyannotation@communion-antiphon\endcsname{%
  \chronologyannotationgroup{composition}{disputed}{\textbf{Composition} -- disputed: \chronologyannotationreach{Matt.5.10}{Matt}{true}\chronologyannotationreach{Matt.5.8}{Matt}{true}\chronologyannotationreach{Matt.5.9}{Matt}{true}\chronologyannotationclaim{composition.gospel-of-matthew}{composition}{catholic-traditional-v1}{disputed}{about A.D. 38-45}{c. A.D. 38--45}; \chronologyannotationreach{Matt.5.10}{Matt}{true}\chronologyannotationreach{Matt.5.8}{Matt}{true}\chronologyannotationreach{Matt.5.9}{Matt}{true}\chronologyannotationclaim{composition.gospel-of-matthew}{composition}{catholic-traditional-v1}{disputed}{about the year 40-42}{c. A.D. 40--42}; \chronologyannotationreach{Matt.5.10}{Matt}{true}\chronologyannotationreach{Matt.5.8}{Matt}{true}\chronologyannotationreach{Matt.5.9}{Matt}{true}\chronologyannotationclaim{composition.gospel-of-matthew}{composition}{catholic-traditional-v1}{disputed}{the years 40-45}{A.D. 40--45}; \chronologyannotationreach{Matt.5.10}{Matt}{true}\chronologyannotationreach{Matt.5.8}{Matt}{true}\chronologyannotationreach{Matt.5.9}{Matt}{true}\chronologyannotationclaim{composition.gospel-of-matthew}{composition}{catholic-traditional-v1}{disputed}{about the year 60-68}{c. A.D. 60--68}; \chronologyannotationreach{Matt.5.10}{Matt}{true}\chronologyannotationreach{Matt.5.8}{Matt}{true}\chronologyannotationreach{Matt.5.9}{Matt}{true}\chronologyannotationclaim{composition.gospel-of-matthew}{composition}{catholic-traditional-v1}{disputed}{about the years 64-67}{c. A.D. 64--67}; \chronologyannotationreach{Matt.5.10}{Matt}{true}\chronologyannotationreach{Matt.5.8}{Matt}{true}\chronologyannotationreach{Matt.5.9}{Matt}{true}\chronologyannotationclaim{composition.gospel-of-matthew}{composition}{catholic-traditional-v1}{disputed}{about the year 50}{c. A.D. 50}.}%
}
% presentation: orientation-v2
\newcommand{\chronologyorientationrow}[2]{%
  \ifcsname triptychchronologyorientation@#1\endcsname
    \csname triptychchronologyorientation@#1\endcsname{#2}%
  \else
    \PackageError{triptych}{No chronology orientation row for #1}{}%
  \fi}
\expandafter\def\csname triptychchronologyorientation@row-f5cbb3bd9936e1a8\endcsname#1{%
Responsorial Psalm: Psalm 24:1bc-2,3--4ab,5--6, Cf. Psalm 24:6 (Hebrew numbering; Vulgate 23) & Historical setting. #1 & Occasion date unresolved\\[.25em]%
}
\expandafter\def\csname triptychchronologyorientation@row-e3f82d7d7c816aeb\endcsname#1{%
Responsorial Psalm: Psalm 24:1bc-2,3--4ab,5--6, Cf. Psalm 24:6 (Hebrew numbering; Vulgate 23) & Composition: Psalms. #1 & \chronologyannotationreach{Ps.23.1}{Ps}{true}\chronologyannotationreach{Ps.23.2}{Ps}{true}\chronologyannotationreach{Ps.23.3}{Ps}{true}\chronologyannotationreach{Ps.23.4}{Ps}{true}\chronologyannotationreach{Ps.23.5}{Ps}{true}\chronologyannotationreach{Ps.23.6}{Ps}{true}\chronologyannotationclaim{critical.psalms.latest-composition-boundary}{composition}{catholic-critical-v1}{preferred}{before the Maccabean period, around 165 B.C.; no individual psalm can be dated securely}{Before c. 165 B.C.}\\[.25em]%
}
\expandafter\def\csname triptychchronologyorientation@row-5a9c87b5fd78fefb\endcsname#1{%
Gospel: Matthew 5:1--12a & Event. #1 & Date unresolved\\[.25em]%
}
\expandafter\def\csname triptychchronologyorientation@row-6661a9700ce1b710\endcsname#1{%
Gospel acclamation: Matthew 11:28; Gospel: Matthew 5:1--12a; Communion antiphon: Matthew 5:8--10 & Composition: The Gospel of St Matthew. #1 & Disputed: \chronologyannotationreach{Matt.11.28}{Matt}{true}\chronologyannotationreach{Matt.5.1}{Matt}{true}\chronologyannotationreach{Matt.5.10}{Matt}{true}\chronologyannotationreach{Matt.5.11}{Matt}{true}\chronologyannotationreach{Matt.5.12}{Matt}{true}\chronologyannotationreach{Matt.5.2}{Matt}{true}\chronologyannotationreach{Matt.5.3}{Matt}{true}\chronologyannotationreach{Matt.5.4}{Matt}{true}\chronologyannotationreach{Matt.5.5}{Matt}{true}\chronologyannotationreach{Matt.5.6}{Matt}{true}\chronologyannotationreach{Matt.5.7}{Matt}{true}\chronologyannotationreach{Matt.5.8}{Matt}{true}\chronologyannotationreach{Matt.5.9}{Matt}{true}\chronologyannotationclaim{composition.gospel-of-matthew}{composition}{catholic-traditional-v1}{disputed}{about A.D. 38-45}{c. A.D. 38--45}; \chronologyannotationreach{Matt.11.28}{Matt}{true}\chronologyannotationreach{Matt.5.1}{Matt}{true}\chronologyannotationreach{Matt.5.10}{Matt}{true}\chronologyannotationreach{Matt.5.11}{Matt}{true}\chronologyannotationreach{Matt.5.12}{Matt}{true}\chronologyannotationreach{Matt.5.2}{Matt}{true}\chronologyannotationreach{Matt.5.3}{Matt}{true}\chronologyannotationreach{Matt.5.4}{Matt}{true}\chronologyannotationreach{Matt.5.5}{Matt}{true}\chronologyannotationreach{Matt.5.6}{Matt}{true}\chronologyannotationreach{Matt.5.7}{Matt}{true}\chronologyannotationreach{Matt.5.8}{Matt}{true}\chronologyannotationreach{Matt.5.9}{Matt}{true}\chronologyannotationclaim{composition.gospel-of-matthew}{composition}{catholic-traditional-v1}{disputed}{about the year 40-42}{c. A.D. 40--42}; \chronologyannotationreach{Matt.11.28}{Matt}{true}\chronologyannotationreach{Matt.5.1}{Matt}{true}\chronologyannotationreach{Matt.5.10}{Matt}{true}\chronologyannotationreach{Matt.5.11}{Matt}{true}\chronologyannotationreach{Matt.5.12}{Matt}{true}\chronologyannotationreach{Matt.5.2}{Matt}{true}\chronologyannotationreach{Matt.5.3}{Matt}{true}\chronologyannotationreach{Matt.5.4}{Matt}{true}\chronologyannotationreach{Matt.5.5}{Matt}{true}\chronologyannotationreach{Matt.5.6}{Matt}{true}\chronologyannotationreach{Matt.5.7}{Matt}{true}\chronologyannotationreach{Matt.5.8}{Matt}{true}\chronologyannotationreach{Matt.5.9}{Matt}{true}\chronologyannotationclaim{composition.gospel-of-matthew}{composition}{catholic-traditional-v1}{disputed}{the years 40-45}{A.D. 40--45}; \chronologyannotationreach{Matt.11.28}{Matt}{true}\chronologyannotationreach{Matt.5.1}{Matt}{true}\chronologyannotationreach{Matt.5.10}{Matt}{true}\chronologyannotationreach{Matt.5.11}{Matt}{true}\chronologyannotationreach{Matt.5.12}{Matt}{true}\chronologyannotationreach{Matt.5.2}{Matt}{true}\chronologyannotationreach{Matt.5.3}{Matt}{true}\chronologyannotationreach{Matt.5.4}{Matt}{true}\chronologyannotationreach{Matt.5.5}{Matt}{true}\chronologyannotationreach{Matt.5.6}{Matt}{true}\chronologyannotationreach{Matt.5.7}{Matt}{true}\chronologyannotationreach{Matt.5.8}{Matt}{true}\chronologyannotationreach{Matt.5.9}{Matt}{true}\chronologyannotationclaim{composition.gospel-of-matthew}{composition}{catholic-traditional-v1}{disputed}{about the year 60-68}{c. A.D. 60--68}; \chronologyannotationreach{Matt.11.28}{Matt}{true}\chronologyannotationreach{Matt.5.1}{Matt}{true}\chronologyannotationreach{Matt.5.10}{Matt}{true}\chronologyannotationreach{Matt.5.11}{Matt}{true}\chronologyannotationreach{Matt.5.12}{Matt}{true}\chronologyannotationreach{Matt.5.2}{Matt}{true}\chronologyannotationreach{Matt.5.3}{Matt}{true}\chronologyannotationreach{Matt.5.4}{Matt}{true}\chronologyannotationreach{Matt.5.5}{Matt}{true}\chronologyannotationreach{Matt.5.6}{Matt}{true}\chronologyannotationreach{Matt.5.7}{Matt}{true}\chronologyannotationreach{Matt.5.8}{Matt}{true}\chronologyannotationreach{Matt.5.9}{Matt}{true}\chronologyannotationclaim{composition.gospel-of-matthew}{composition}{catholic-traditional-v1}{disputed}{about the years 64-67}{c. A.D. 64--67}; \chronologyannotationreach{Matt.11.28}{Matt}{true}\chronologyannotationreach{Matt.5.1}{Matt}{true}\chronologyannotationreach{Matt.5.10}{Matt}{true}\chronologyannotationreach{Matt.5.11}{Matt}{true}\chronologyannotationreach{Matt.5.12}{Matt}{true}\chronologyannotationreach{Matt.5.2}{Matt}{true}\chronologyannotationreach{Matt.5.3}{Matt}{true}\chronologyannotationreach{Matt.5.4}{Matt}{true}\chronologyannotationreach{Matt.5.5}{Matt}{true}\chronologyannotationreach{Matt.5.6}{Matt}{true}\chronologyannotationreach{Matt.5.7}{Matt}{true}\chronologyannotationreach{Matt.5.8}{Matt}{true}\chronologyannotationreach{Matt.5.9}{Matt}{true}\chronologyannotationclaim{composition.gospel-of-matthew}{composition}{catholic-traditional-v1}{disputed}{about the year 50}{c. A.D. 50}\\[.25em]%
}
\expandafter\def\csname triptychchronologyorientation@row-4f310c075dd990dc\endcsname#1{%
Second Reading: 1 John 3:1--3 & Composition: The First Epistle of St John. #1 & \chronologyannotationreach{1John.3.1}{1John}{true}\chronologyannotationreach{1John.3.2}{1John}{true}\chronologyannotationreach{1John.3.3}{1John}{true}\chronologyannotationclaim{composition.first-epistle-of-st-john}{composition}{catholic-traditional-v1}{preferred}{from the year 90 to 100 (approximately)}{c. A.D. 90--100}\\[.25em]%
}
\expandafter\def\csname triptychchronologyorientation@row-02f156d681e4660a\endcsname#1{%
First Reading: Revelation 7:2--4,9--14 & Composition: The Apocalypse of St John. #1 & Disputed: \chronologyannotationreach{Apoc.7.10}{Apoc}{true}\chronologyannotationreach{Apoc.7.11}{Apoc}{true}\chronologyannotationreach{Apoc.7.12}{Apoc}{true}\chronologyannotationreach{Apoc.7.13}{Apoc}{true}\chronologyannotationreach{Apoc.7.14}{Apoc}{true}\chronologyannotationreach{Apoc.7.2}{Apoc}{true}\chronologyannotationreach{Apoc.7.3}{Apoc}{true}\chronologyannotationreach{Apoc.7.4}{Apoc}{true}\chronologyannotationreach{Apoc.7.9}{Apoc}{true}\chronologyannotationclaim{composition.apocalypse}{composition}{catholic-traditional-v1}{disputed}{the reign of Domitian ... A.D. 81-96}{A.D. 81--96}; \chronologyannotationreach{Apoc.7.10}{Apoc}{true}\chronologyannotationreach{Apoc.7.11}{Apoc}{true}\chronologyannotationreach{Apoc.7.12}{Apoc}{true}\chronologyannotationreach{Apoc.7.13}{Apoc}{true}\chronologyannotationreach{Apoc.7.14}{Apoc}{true}\chronologyannotationreach{Apoc.7.2}{Apoc}{true}\chronologyannotationreach{Apoc.7.3}{Apoc}{true}\chronologyannotationreach{Apoc.7.4}{Apoc}{true}\chronologyannotationreach{Apoc.7.9}{Apoc}{true}\chronologyannotationclaim{composition.apocalypse}{composition}{catholic-traditional-v1}{disputed}{from the year 90 to 100 (approximately)}{c. A.D. 90--100}\\[.25em]%
}

\begin{chronologyorientationtable}
\chronologyorientationrow{row-f5cbb3bd9936e1a8}{Seeking admission to the Lord's mountain and sanctuary; original occasion unknown.}
\chronologyorientationrow{row-e3f82d7d7c816aeb}{Individual writing place unknown. The date is the collection's latest bound, not an individual composition year.}
\chronologyorientationrow{row-5a9c87b5fd78fefb}{Jesus teaches disciples on a mountain during his Galilean ministry; exact hill unknown (Matthew 4:23--5:2).}
\chronologyorientationrow{row-6661a9700ce1b710}{Tradition associates an earlier Hebrew form with Palestine; the Greek text's relation is disputed. These estimates concern differing accounts of composition.}
\chronologyorientationrow{row-4f310c075dd990dc}{Probably associated with Ephesus; recipients' town uncertain. The approximate span concerns the Johannine writings collectively.}
\chronologyorientationrow{row-02f156d681e4660a}{The seer names Patmos; the book addresses seven churches in Asia. The throne in chapter 7 belongs to the vision.}
\end{chronologyorientationtable}
\smallskip
\noindent\textit{Historical evidence.} The expansive study's Appendix: Scriptural Date and Location gives the complete explanations under ``Psalm 24: a sanctuary setting and a collection bound,'' ``Matthew: the mountain and the writing of the Gospel,'' ``First John: probable Ephesian setting,'' and ``Revelation: Patmos and the churches in Asia.'' Its sources are the USCCB Psalms introduction; Jacquier on Matthew; Drum and Durand on First John; and Van den Biesen and Durand on Revelation. The older encyclopedia articles report historical scholarly judgments and their alternatives.
\zlabel{triptych:concise:chronology:end}
